\documentclass[../../../include/open-logic-section]{subfiles}

\begin{document}

\olfileid{sth}{z}{pairs}
\olsection{الأزواج}

المسلّمة التالية التي ينبغي النظر فيها هي:

\begin{axiom}[الأزواج]
لكل مجموعتين~$a,b$، توجد المجموعة~$\{a,b\}$.
\[
	\forall a \forall b \exists P \forall x (x \in P \liff (x = a \lor x = b))
\]
\end{axiom}

وهكذا نبرر هذه المسلّمة باستعمال التصور التكراري. افترض أن~$a$ متاحة
في المرحلة~$S$، وأن~$b$ متاحة في المرحلة~$T$. ولتكن~$M$ المرحلة
المتأخرة من~$S$ و$T$. ولما كانت~$a$ و$b$ كلتاهما متاحتين في
المرحلة~$M$، كانت المجموعة~$\{a,b\}$ تجمعًا ممكنًا متاحًا في أي مرحلة
تأتي بعد~$M$.

لكن تمهل!{} لماذا نفترض وجود أي مراحل بعد~$M$؟ إذا لم توجد، فشل
تبريرنا. ولذلك، لكي نبرر مسلّمة الأزواج، يلزمنا أن نضيف مبدأ آخر إلى
الرواية التي قدمناها في~\olref[sth][z][story]{sec}، وهو:
\begin{enumerate}
	\item[] \stagessucc. لا توجد مرحلة أخيرة.
\end{enumerate}
فهل هذا المبدأ مبرر؟ لم يصرح شيء في رواية شونفيلد \emph{صراحة} بأنه
لا توجد مرحلة أخيرة. ومع ذلك، فحتى لو كان (على وجه الدقة) إضافة إلى
روايتنا، فإنه ينسجم جيدًا مع الفكرة الأساسية القائلة إن المجموعات
تتكون في مراحل. وسنقبله ببساطة فيما يأتي، ولذلك سنقبل مسلّمة الأزواج
أيضًا.

وباستخدام هذه المسلّمة الجديدة نستطيع إثبات وجود مجموعات أخرى كثيرة.
مثلًا:

\begin{prop}\ollabel{prop:pairsconsequences}
لكل مجموعتين~$a$ و$b$، توجد المجموعات الآتية:
	\begin{enumerate}
		\item\ollabel{singleton} $\{a\}$
		\item\ollabel{binunion} $a \cup b$
		\item\ollabel{tuples} $\tuple{a, b}$
	\end{enumerate}
\end{prop}

\begin{proof}
\olref{singleton}. بموجب الأزواج، توجد~$\{a, a\}$، وهي~$\{a\}$
بموجب الامتدادية.

\olref{binunion}. بموجب الأزواج، توجد~$\{a, b\}$. وتوجد الآن
$a \cup b = \bigcup \{a, b\}$ بموجب الاتحاد.

\olref{tuples}. بموجب~\olref{singleton}، توجد~$\{a\}$. وبموجب الأزواج،
توجد~$\{a,b\}$. وتوجد الآن
$\{\{a\}, \{a, b\}\} = \tuple{a, b}$ بتطبيق الأزواج مرة أخرى.
\end{proof}

\begin{prob}
بيّن أنه لكل مجموعات~$a,b,c$، توجد المجموعة~$\{a,b,c\}$.
\end{prob}

\begin{prob}
بيّن أنه لكل مجموعات~$a_1,\ldots,a_n$، توجد المجموعة
$\{a_1,\ldots,a_n\}$.
\end{prob}

%\begin{proof} بتطبيق الأزواج مرتين، توجد $\{\{a_1, a_2\}, \{a_1,
%   a_3\}\}$. وبموجب الاتحاد والامتدادية، توجد $\bigcup \{\{a_1,
%   a_2\}, \{a_1, a_3\}\} = \{a_1, a_2, a_3\}$. كرر هذه
%   الحيلة كلما لزم. \end{proof}

\end{document}
